\documentclass{ximera}
\title{Fun with Vectors}

\newcommand{\pskip}{\vskip 0.1 in}

\begin{document}
\begin{abstract}
%Vector basics.
\end{abstract}
\maketitle

\section{Relative Position}

\begin{example} \label{Ex:dlfkjef3}
For the following questions draw pictures with the non-colinear points $P$, $Q$, and $R$ along with the relevant vectors to help with your explanations.

\begin{enumerate}
\item What vectors describe the position of $P$ relative to $Q$ and the position of $Q$ relative to $R$?

\item What vector describes the position of $R$ relative to $P$?

\item Express the third vector above in terms of the first two.

\item Express the vector giving the position of $R$ relative to $Q$ in terms of the vectors giving the positions of $R$ relative to $P$ and of $Q$ relative to $P$. 
\end{enumerate}
\end{example}

\begin{example} \label{Exlfl3ef333}
Let $A$, $B$, and $C$ be non-colinear points.

\begin{enumerate}
\item Draw the vectors $\overrightarrow{AB}$, $\overrightarrow{AC}$, and $\overrightarrow{BC}$.

\item Express $\overrightarrow{BC}$ in terms of $\overrightarrow{AB}$ and $\overrightarrow{AC}$.

\item Explain the meaning of each of these vectors in terms of relative position.

\item Draw the vector 
\[
 \overrightarrow{AB} + \overrightarrow{BC} + \overrightarrow{CA}
\]
and explain its meaning in terms of relative position.

\item True or false:
\[
    \left|  \overrightarrow{AB} + \overrightarrow{BC} \right| =  \left|  \overrightarrow{AB}\right| + \left| \overrightarrow{BC} \right|
\]
Explain. Draw a picture to help.
\end{enumerate}
\end{example}

\begin{example} \label{Exkdk32k21}
Points $A$ and $B$ have respective coordinates $(3,1)$ and $(-2,4)$. Point $O$ is the origin. Point $C$ lies on the line $AB$ and is $10$ units from $A$.


\begin{enumerate}
\item Ignore the given information and express $\overrightarrow{OC}$ in terms of $\overrightarrow{OA}$ and $\overrightarrow{AC}$ for any three points $O$, $A$, and $C$.

\item Find the components of the vector $\overrightarrow{AB}$.

\item Find the components of the unit vector (a vector having length one) in the direction of $\overrightarrow{AB}$.

\item Find the components of the vector $\overrightarrow{AC}$. Give all possibilities.

\item Find the coordinates of $C$. Give all possibilities.
\end{enumerate}
\end{example}

\begin{example} \label{Exf3r3rnms}
Quadrilateral $ABCD$ is a parallelogram (opposite sides are parallel and therefore have equal lengths). Point $O$ is the origin.

\begin{enumerate}
\item Write an equation relating the vectors  $\overrightarrow{AB}$ and $\overrightarrow{CD}$.

\item Express the vector $\overrightarrow{OD}$ in terms of the vectors $\overrightarrow{OA}$, $\overrightarrow{OB}$, and $\overrightarrow{OC}$.

\item Suppose $A$, $B$, and $C$ have respective coordinates $(2,1)$, $(5,8)$, and $(11,14)$. Determine the coordinates of $D$. 
\end{enumerate}
\end{example}


\begin{example} \label{ex8df83r3r}
Let $O$ be the origin and $P$ the point with coordinates $(5,2)$. Let $\theta$, $0 < \theta < \pi/2$, be the polar angle of $\overrightarrow{OP}$. Let $OPQR$ be the counterclockwise square on $\overline{OP}$

\begin{enumerate}
\item Express $\overrightarrow{OQ}$ in terms of $\overrightarrow{OP}$ and $\overrightarrow{OR}$.

\item Find the components of the vectors $\overrightarrow{OP}$ and $\overrightarrow{OR}$. Use these to find the components of the vector $\overrightarrow{OQ}$.

\item Use the result of part (b) to find the exact values of $\cos(\theta + \pi/4)$ and $\sin(\theta + \pi/4)$.

\item Solve this problem in general to express $\cos(\theta + \pi/4)$ and $\sin(\theta + \pi/4)$ in terms of $\cos\theta$ and $\sin\theta$ for any angle $\theta$.

\end{enumerate}
\end{example}

\begin{example} \label{ex8eef83r3r}
Let $O$ be the origin and $P$ the point with coordinates $(5,2)$. Let $\theta$, $0 < \theta < \pi/2$, be the polar angle of $\overrightarrow{OP}$. Let $OPQ$ be the counterclockwise equilateral triangle on $\overline{OP}$ and let $M$ be the midpoint of $\overline{OP}$.

\begin{enumerate}
\item Express $\overrightarrow{OQ}$ in terms of $\overrightarrow{OM}$ and $\overrightarrow{MQ}$.

\item Find the components of the vectors $\overrightarrow{OM}$ and $\overrightarrow{MQ}$. Use these to find the components of the vector $\overrightarrow{OQ}$.

\item Use the result of part (b) to find the exact values of $\cos(\theta + \pi/3)$ and $\sin(\theta + \pi/3)$.

\item Solve this problem in general to express $\cos(\theta + \pi/3)$ and $\sin(\theta + \pi/3)$ in terms of $\cos\theta$ and $\sin\theta$ for any angle $\theta$.

\item Use a similar method to express $\cos(\theta - \pi/3)$ and $\sin(\theta - \pi/3)$ in terms of $\cos\theta$ and $\sin\theta$ for any angle $\theta$.

\end{enumerate}
\end{example}








\end{document}